%==============================================================================
% Capitulo 3: Estado del Arte
%==============================================================================
\section{Estado del arte}

Este capítulo revisa la literatura académica y las soluciones comerciales disponibles en tres áreas relevantes para el presente trabajo: plataformas de mesa de ayuda, herramientas de orquestación de flujos, y métodos de clasificación automática de tickets de soporte. El proposito es situar la contribucion de esta investigación en el contexto de los desarrollos previos y fundamentar la brecha que motiva el trabajo.

\subsection{Soluciones comerciales de mesa de ayuda}

El segmento de soluciones comerciales para gestión de mesa de ayuda se encuentra dominado por plataformas integrales que cubren el ciclo completó de vida del ticket. Zendesk Suite ofrece funcionalidades de recepción multicanal, automatización de flujos básicos y, mediante un modulo adicional, capacidades de inteligencia artificial para sugerencia de respuestas y defleccion de consultas frecuentes. Freshdesk propone una arquitectura comparable e incorpora el motor Freddy AI para clasificación automática y enrutamiento inteligente. Jira Service Management, parte del ecosistema de Atlassian, ofrece capacidades limitadas de automatización inteligente en su plan estandar; la edicion Data Center permite el despliegue autoalojado para organizaciones con requisitos de soberanía sobre los datos. ServiceNow ITSM constituye la opcion de referencia en el segmento corporativo, con un modulo Predictive Intelligence que aplica aprendizaje automático sobre el histórico de tickets de la organización.

La adopción de capacidades de inteligencia artificial en plataformas comerciales de mesa de ayuda ha crecido de manera sostenida en los últimos anos, particularmente en organizaciones medianas de mercados de habla inglesa. No obstante, estas plataformas presentan dos limitaciones relevantes para el contexto de la presente investigación. La primera es el modelo de costos basado en agente y mes, que escala linealmente con el tamaño del equipo de soporte. La segunda es que la mayoría de estas plataformas operan exclusivamente en modalidad de software como servicio (SaaS) con almacenamiento de datos en jurisdicciones extranjeras, lo que introduce consideraciones adicionales sobre transferencia internacional de datos personales bajo la Ley~25.326 \parencite{Ley25326_2000}.

En el segmento de código abierto, osTicket constituye una alternativa de despliegue autoalojado sin costo de licencia, aunque carece de capacidades nativas de clasificación automática mediante procesamiento de lenguaje natural moderno. Esta brecha funcional motiva la busqueda de arquitecturas compositivas que integren componentes maduros de código abierto con servicios de inferencia bajo demanda.

\subsection{Plataformas de orquestacion de flujos}

Dentro del segmento de plataformas de orquestación de flujos, las soluciones mas difundidas incluyen Zapier, Make (anteriormente Integromat), Microsoft Power Automate, Apache Airflow y N8N. Las primeras tres operan exclusivamente como SaaS, lo que las descarta como base para implementaciones donde el control sobre la infraestructura es un requisito explicito. Apache Airflow se orienta a flujos de trabajo programados de tipo procesamiento de datos por lotes, con un diseñó centrado en grafos aciclicos dirigidos (DAG) definidos como código Python \parencite{Apache2024_airflow}; aunque potente, su paradigma de ejecución programada resulta inadecuado para un caso de uso reactivo y orientado a eventos como el de una mesa de ayuda.

N8N ofrece un modelo de construcción visual de flujos sobre nodos predefinidos, soporte nativo para webhooks y disparadores reactivos, ejecución basada en eventos y posibilidad de despliegue completamente autoalojado bajo Docker \parencite{n8n2024_docs}. La disponibilidad de su código fuente bajo licencia Sustainable Use facilita la auditoria de seguridad y la adaptacion a requerimientos específicos sin generar dependencia exclusiva del proveedor. Estas caracteristicas confluyen para posicionar a N8N como la opcion mas adecuada para el escenario propuesto, donde la soberanía sobre los datos y los flujos de procesamiento constituye un requisito no negociable.

\subsection{Trabajos académicos previos sobre clasificación de tickets}

La literatura académica sobre clasificación automática de tickets de soporte ha experimentado un crecimiento considerable durante los últimos anos. A continuación se presentan los trabajos mas relevantes, organizados por enfoque metodológico y por orden cronologico.

\textcite{Paramesh2019_it_service_desk} constituyen uno de los primeros trabajos exhaustivos en el área. Utilizando combinaciones de maquinas de vectores de soporte (SVM) y representaciones TF-IDF sobre un corpus en inglés de aproximadamente 10.000 tickets, reportan exactitudes cercanas al 87~\%. Su contribucion principal radica en la demostracion de que las tecnicas clasicas de aprendizaje supervisado, cuando se aplican sobre representaciones textuales adecuadas, pueden alcanzar niveles de desempeño operativamente utiles para la clasificación automática de tickets.

\textcite{Revina2020_bert_tickets} extienden este enfoque incorporando arquitecturas basadas en BERT y reportan mejoras significativas, alcanzando exactitudes del orden del 91~\% sobre un corpus de tamaño comparable. Su hallazgo mas relevante para el presente trabajo es la constatacion de que modelos preentrenados sobre grandes corpus de texto general pueden transferir conocimiento linguistico útil al dominio específico de los tickets de soporte, incluso sin ajuste fino sobre datos del dominio. Este resultado respalda la estrategia adoptada en la presente investigación de utilizar un LLM de proposito general como motor de inferencia.

La literatura reciente sobre evaluación de modelos de lenguaje en escenarios de pocos ejemplos ha documentado que el rendimiento de clasificación depende fuertemente del idioma del corpus, con resultados que favorecen a los idiomas mejor representados en los corpus de entrenamiento de los modelos contemporaneos. El espanol, y en particular sus variantes regionales como el rioplatense, ha recibido una atención notablemente menor en las evaluaciones sistematicas publicadas. Esta brecha constituye una de las motivaciones academicas del presente trabajo, el cual aporta evidencia empírica sobre el desempeño del clasificador en espanol rioplatense aplicado a un dominio de mesa de ayuda de organización mediana.

\subsection{Analisis comparativo y brecha identificada}

La Tabla~\ref{tab:comparativa} presenta la comparación sistematica de las principales soluciones revisadas frente a la propuesta del presente trabajo, considerando cuatro dimensiones relevantes: tipo de despliegue, capacidad de clasificación automática, soberanía sobre los datos y modelo de costos.

\begin{table}[htbp]
\centering
\caption{Comparacion de soluciones para gestion automatizada de mesa de ayuda}
\label{tab:comparativa}
\small
\begin{tabular}{@{}p{2.6cm} p{2.8cm} >{\raggedright\arraybackslash}p{3.5cm} >{\raggedright\arraybackslash}p{3.8cm}@{}}
\toprule
\textbf{Solucion} & \textbf{Despliegue} & \textbf{Clasificacion automatica} & \textbf{Modelo de costo} \\
\midrule
Zendesk Suite     & SaaS comercial      & Si, modulo de IA                & Alto, por agente/mes \\
Freshdesk         & SaaS comercial      & Si, motor Freddy AI             & Medio-alto, por agente \\
Jira Service Mgmt & SaaS o autoalojado  & Limitada en plan estandar       & Medio, por agente \\
ServiceNow ITSM   & SaaS comercial      & Si, Predictive Intelligence     & Muy alto, contrato anual \\
osTicket          & Codigo abierto      & No nativa                       & Sin costo \\
\addlinespace
\textbf{Propuesta} & Hibrida (iPaaS + modulo propio) & Si, LLM + reglas deterministas & Bajo, costo por uso de API \\
\bottomrule
\end{tabular}
\end{table}

La sintesis comparativa evidencia que ninguna de las soluciones revisadas combina simultaneamente las cuatro propiedades que el contexto del trabajo demanda: despliegue autoalojado para resguardo de datos sensibles, costo bajo y predecible en el tiempo, capacidad nativa de clasificación automática mediante NLP moderno, y soporte específico para canales multiples incluyendo telefonia con transcripción automática. La brecha identificada se sintetiza en la ausencia de soluciones de bajo costo, autoalojadas, con clasificación basada en modelos de lenguaje grandes y validadas empíricamente sobre corpus en espanol rioplatense para el dominio específico de mesas de ayuda empresariales medianas.

La presente investigación se ubica precisamente en esa interseccion y propone una arquitectura compositiva que reutiliza componentes maduros de código abierto y servicios de inferencia bajo demanda, manteniendo bajo control de la organización los datos sensibles y los flujos de orquestación, y aportando evidencia empírica reproducible sobre el desempeño del enfoque en un contexto linguistico y organizacional escasamente documentado en la literatura.

\newpage
